\documentclass[11pt]{article}
\usepackage[margin=1in]{geometry}
\usepackage{amsmath,amssymb,amsthm}
\usepackage{mathtools}
\usepackage{enumitem}
\usepackage{booktabs}
\usepackage{hyperref}
\usepackage{xcolor}
\hypersetup{colorlinks=true,linkcolor=blue,citecolor=blue,urlcolor=blue}

\newcommand{\Quant}{\textsf{Quant}}
\newcommand{\Device}{\textsf{Device}}
\newcommand{\Ticket}{\textsf{Ticket}}
\newcommand{\Extract}{\textsf{Extract}}
\newcommand{\Rows}{\textsf{Rows}}
\newcommand{\Hb}{\mathsf{H}}
\newcommand{\MatMul}{\textsf{MatMul}}
\newcommand{\NoisyQuantize}{\textsf{NoisyQuantize}}
\newcommand{\SampleLine}{\textsf{SampleLine}}
\newcommand{\Trunc}{\textsf{Trunc}}
\newcommand{\RZ}{\textsf{RZ}}
\newcommand{\RNE}{\textsf{RNE}}
\newcommand{\ulp}{\textsf{ulp}}
\newcommand{\ufp}{\textsf{ufp}}
\newcommand{\eps}{\varepsilon}
\newcommand{\sg}{\mathrm{sgn}}
\newcommand{\Ct}{\widetilde{C}}
\newcommand{\At}{\widetilde{A}}
\newcommand{\Bt}{\widetilde{B}}
\newcommand{\T}{\mathsf{T}}
\newcommand{\R}{\mathbb{R}}
\newcommand{\FP}[1]{\textsf{FP#1}}
\newcommand{\sigmin}{\sigma_{\min}}

\newtheorem{assumption}{Assumption}
\newtheorem{conjecture}{Conjecture}
\newtheorem{definition}{Definition}
\newtheorem{remark}{Remark}

\title{\textbf{Pearl Floating-Point Scheme Specification:\\ FP16 Accumulation-Hardness Variant}}
\author{Jon Pry\\ \texttt{jonpry@gmail.com}}
\date{October 2026}

\begin{document}
\maketitle

\begin{abstract}
This document specifies an alternative instantiation of the Pearl proof-of-useful-work
protocol in which the unit of mining work is \FP{16} matrix multiplication on GPUs, with
NVIDIA A100 (GA100) tensor cores as the first whitelisted device. It extends the \FP{8}
certificate-v4 specification rather than replacing it: the commitment structure, seed chain,
state window, ticket, target, mixture-of-experts extension, and zero-knowledge split are
unchanged, and the \FP{8} scheme remains available as a separate \Quant{} value. The new
content is a \emph{hardness source that does not rely on coarse quantization}. Where the
\FP{8} scheme makes the lottery product hard to shortcut by rounding onto a coarse grid, the
\FP{16} scheme derives hardness from the \emph{nonlinearity of the device's own accumulation}:
the A100 tensor core truncates each product onto a per-group alignment grid before summing,
an operation that does not commute with the reduction and therefore cannot be expressed as a
matrix multiplication. We give the bit-exact A100 accumulation model (validated against
silicon), a policy check---\emph{unpredictable accumulation steps}---that the verifier
evaluates during its tile replay, and an analysis of the attack surface. The practical payoff
is that recovered products carry \FP{16}-level accuracy rather than \FP{8}-residual accuracy,
at the cost of a new, device-specific hardness assumption. We find \FP{16} robust and
\textsf{BF16} too close to linear to admit the same argument.
\end{abstract}

\tableofcontents

\section{Introduction}\label{sec:intro}

Proof-of-useful-work (PoUW) replaces the artificial puzzle of classical proof of work with a
computation someone already wants performed. The Pearl protocol instantiates the PoUW of
Komargodski and Weinstein~\cite{KW25} for matrix multiplication: a miner multiplies matrices
of its own choosing and every tile of the product doubles as a lottery ticket. The obstacle
is permissionlessness---the miner chooses its inputs, so the protocol must guarantee that no
choice of inputs, however degenerate, makes valid tickets cheaper than the prescribed amount
of work.

The original integer protocol~\cite{KW25} achieves this by hashing the \emph{transcript} of
the computation: the product is formed in $r\times r$ blocks and a random oracle is applied to
all $(n/r)^3$ intermediate blocks, so that even the all-zeroes product is as hard as a generic
one. The Pearl \FP{8} specification~\cite{PearlFP8} instead hashes only output tiles and
recovers hardness from quantization: because entrywise quantization is nonlinear, the low-rank
noise injected before quantization destroys the low-rank shortcut, and tiles can be hashed
directly. This is cheaper than transcript hashing, but its security depends on the rounding
being coarse: the quantization error must be large enough that its contribution to the output
survives the accumulator. At \FP{8} this holds comfortably; as the target format becomes more
precise, the quantization error shrinks and the margin vanishes.

This document describes a third point in the design space, aimed at \FP{16}. We keep output-
tile hashing---so the miner runs its ordinary GEMM with no in-kernel hashing---but we obtain
hardness from a different nonlinearity. Modern tensor cores do not accumulate in exact
arithmetic. On the A100, each product is truncated toward zero onto an alignment grid shared by
a group of eight products, the group is summed, and the running accumulator is rounded after
every group. The truncation is applied \emph{per product, before the reduction}. This is the
same structural device that makes transcript hashing hard---a per-summand operation that a
single contraction cannot reproduce---except that it is performed by the hardware for free, on
every multiply-accumulate the miner already issues. An honest miner's output therefore encodes,
at no extra cost, a quantity that an adversary with the low-rank shortcut cannot cheaply
reconstruct.

\paragraph{Organization.} Section~\ref{sec:abstract} recalls the abstract protocol and states
the new hardness premise. Section~\ref{sec:prelim} fixes notation and the work model.
Section~\ref{sec:device} gives the bit-exact A100 accumulation model---the technical core of
this variant. Section~\ref{sec:hardness} states the hardness argument and assumption.
Section~\ref{sec:policy} specifies the jackpot policy, including the unpredictable-accumulation-
steps check. Section~\ref{sec:attacks} surveys attack vectors. Section~\ref{sec:miner}
describes the honest miner. Section~\ref{sec:compare} compares \FP{16}, \FP{8} and \textsf{BF16}.
The appendices give the A100 arithmetic specification, the \FP{16} commit type, the noise and
quantization procedures, the certified-work ratio, and a summary of the empirical validation.

\section{The abstract protocol}\label{sec:abstract}

Let $A\in\R^{m\times k}$ and $B\in\R^{n\times k}$, with useful product $C:=AB^\T$. As in the
\FP{8} scheme, the protocol derives seed-dependent rank-$r$ noise matrices $N_A\in\R^{m\times k}$
and $N_B\in\R^{n\times k}$ and computes
\begin{equation}
\Ct := Q(A+N_A)\,Q(B+N_B)^\T,
\end{equation}
where $Q$ is the protocol's quantization into the committed storage format and the
multiplication is carried out by the committed device's native kernel. Each noise matrix is
unpredictable before the inputs are committed. The product $\Ct$ drives the lottery; the honest
miner recovers the useful product by the cheap low-rank correction
$\widehat C := \Ct - A N_B^\T - N_A B^\T - N_A N_B^\T$, exactly as in~\cite{KW25,PearlFP8}.

The difference from the \FP{8} scheme is in \emph{what makes $\Ct$ hard to shortcut}. Setting
$A=B=0$ isolates the lottery product as $Q(N_A)\,Q(N_B)^\T$. The \FP{8} premise is that
entrywise quantization raises the rank of $Q(N_A)$ enough that this product is as hard as a
generic quantized product. The \FP{16} premise is weaker on quantization and stronger on the
device: at \FP{16} the quantization barely raises the rank, but the device does not compute the
exact product $Q(N_A)Q(N_B)^\T$---it computes a \emph{per-product-truncated} version of it,
tile by tile, and it is that version the ticket hashes.

\paragraph{The central object.} For operands stored in a format with $p$-bit significands, write
the device dot product of row $i$ of $\At:=Q(A+N_A)$ with row $j$ of $\Bt:=Q(B+N_B)$ as
\begin{equation}\label{eq:dotdevice}
D_{ij} \;=\; \textsf{Acc}\big(\At_{i,0}\Bt_{j,0},\,\dots,\,\At_{i,k-1}\Bt_{j,k-1}\big),
\end{equation}
where $\textsf{Acc}$ is the device accumulation of Section~\ref{sec:device}, \emph{not} the
exact sum. The exact sum $S_{ij}=\sum_u \At_{iu}\Bt_{ju}$ is, for low-rank-plus-noise operands,
computable by an adversary in $O(r)$ work per cell through the low-rank structure. The quantity
$D_{ij}$ differs from $\RZ(S_{ij})$ by the accumulated per-product truncation remainders, and
reconstructing those remainders is what the protocol charges for. Section~\ref{sec:hardness}
makes this precise.

\section{Preliminaries}\label{sec:prelim}

\paragraph{Notation.} Matrices are over the reals, stored in a committed floating-point format.
$C:=AB^\T\in\R^{m\times n}$; $C_{ij}$ is the inner product of row $i$ of $A$ with row $j$ of $B$.
For selected row-index sets $I_A,I_B$ (chosen by periodic partitions $\mathcal P_A,\mathcal P_B$,
as in the \FP{8} scheme), $\Ct_{I_A,I_B}:=\At_{I_A}(\Bt_{I_B})^\T$ is the recomputed tile. The
seeds $\texttt{noise\_seed}_A,\texttt{noise\_seed}_B$ determine the noise. The ticket digest is a
keyed function $\Ticket$ from tiles to $\{0,1\}^\lambda$, $\lambda=256$:
\[
\Ticket(\Ct_{I_A,I_B}) := \Hb_{\text{``jackpot''}}\big(\Extract(\Ct_{I_A,I_B},\mathcal P_A,\mathcal P_B);\,\texttt{noise\_seed}_A\big).
\]
All of $\Ticket$, $\Extract$, the seed chain, the state window of depth $D$, the domain-separated
BLAKE3 hashing, and the commitment format are shared with the \FP{8} specification~\cite{PearlFP8}
and are not restated here except where \FP{16} changes them (Appendix~\ref{app:commit}).

\paragraph{Work model.} Work is measured in multiply-accumulate operations (MACs). The direct
cost of the selected tile is $T_{\text{tile}}:=|I_A|\cdot|I_B|\cdot k$. As in~\cite{PearlFP8} the
MAC count fixes scale only; it excludes quantization, noise generation, hashing, and memory
movement. The novelty of this variant is that the quantity securing the tile is produced
\emph{inside} these MACs by the device, so no work beyond the GEMM is charged to the miner.

\paragraph{Device and Quant.} $\Device$ and $\Quant$ are consensus-allowed enumerations. This
document adds $\Device=\textsf{A100}$ with the accumulation model of Section~\ref{sec:device}, and
a $\Quant$ value whose storage type is \FP{16} and whose commit type is the \FP{16} block format
of Appendix~\ref{app:commit}. The \FP{8} values remain available unchanged.

\section{The A100 accumulation model}\label{sec:device}

We specify the arithmetic of the A100 \texttt{HMMA.16816.F32} instruction (and the
equivalent \texttt{m16n8k8}/\texttt{m16n8k16} and \texttt{wmma} forms, which chain identically)
for \FP{16} and \textsf{BF16} inputs with \FP{32} accumulation. The model is validated bit-exact
against silicon in Appendix~\ref{app:validation}. The verifier's implementation is authoritative;
the description below is phrased in exact arithmetic.

Let the $k$-axis be partitioned into \emph{groups} of $G=8$ consecutive products, processed in
order (two groups per \texttt{k16} instruction). The internal accumulator precision is
$W=24$ significant bits (full \FP{32} significand). For a stored value $x$, let $\eps(x)$ be its
unbiased stored exponent, with subnormal inputs taking the format minimum ($-14$ for \FP{16},
$-126$ for \textsf{BF16}). For a nonzero \FP{32} accumulator $c$, let $\eps(c):=\max(\lfloor\log_2|c|\rfloor,-126)$.
Let $\Trunc(x,g)$ denote truncation of $x$ toward zero onto the grid $2^{g}$ (sign-magnitude).

\paragraph{One group step.} With incoming accumulator $c$ and group products $p_u=\At_{iu}\Bt_{ju}$
(each exact; its exponent is $\eps(\At_{iu})+\eps(\Bt_{ju})$):
\begin{enumerate}[nosep]
\item If the group has no nonzero product and $c$ is unchanged, the step is a no-op.
\item Compute the alignment exponent
      $\eta := \max\big(\{\eps(\At_{iu})+\eps(\Bt_{ju}) : p_u\neq0\}\cup\{\eps(c)\}\big)$,
      where $c=0$ does not participate and a subnormal $c$ contributes $-126$.
\item Truncate every product and the accumulator onto the grid $2^{\eta-W}$:
      $\bar p_u := \Trunc(p_u,\eta-W)$, $\bar c := \Trunc(c,\eta-W)$.
\item Sum the integers exactly: $S := \bar c + \sum_u \bar p_u$.
\item Round $S$ toward zero to \FP{32}: $c' := \RZ_{\FP{32}}(S)$, i.e.\ keep the top $24$
      significant bits, place subnormal results on the $2^{-149}$ grid, and map $|S|\ge 2^{128}$
      to $\pm\infty$.
\end{enumerate}
The final accumulator after all groups is $D_{ij}$. Boundary behavior: a zero result is always
$+0$; $\infty$ and NaN follow IEEE-754 on the inputs (NaN sign and payload are unspecified); an
overflowed accumulator is \emph{sticky}---a later opposite-sign overflow does not produce NaN,
and the first overflow's sign is retained. The verifier rejects any tile whose operands or
intermediates are non-finite (Appendix~\ref{app:arith}).

\paragraph{The key feature.} Step 3 truncates each product \emph{individually} before step 4
sums them. Write $p'$ for the product significand width (the product of two $p$-bit significands
carries $p'=2p$ significant bits: $p'=22$ for \FP{16}, $p'=16$ for \textsf{BF16}). Once the running
accumulator exceeds a group's products by more than about $W-p'$ bits---roughly a factor
$2^{W-p'}$---the low bits of each product fall below the grid $2^{\eta-W}$ and are discarded. For
\FP{16}, $p'=22$ and $W=24$, so a product loses bits as soon as the accumulator is more than
$\sim\!4\times$ ($2^{W-p'}=2^2$) its magnitude; for a realistic dot product this happens for the
majority of products. For \textsf{BF16}, $p'=16$, so products survive the window unless they are
more than $W-p'=8$ binades below $\eta$; far fewer are truncated (Section~\ref{sec:compare}).

\paragraph{Breakpoints.} Call a group step a \emph{breakpoint} if the accumulator truncation in
step 3 or the round in step 5 discards a nonzero bit, i.e.\ if $\bar c\neq c$ or
$\RZ_{\FP{32}}(S)\neq S$. Per-product truncations ($\bar p_u\neq p_u$) are \emph{not} breakpoints;
they are counted separately as $N_{\text{pt}}$ (Section~\ref{sec:policy}), which is why $N_{\text{pt}}$
is well defined as the truncated products in non-breakpoint steps. The density of breakpoints over a
tile, $f_{\text{bp}}$, measures how much of the accumulation is irrecoverable from the exact sum
alone; it is the quantity the policy of Section~\ref{sec:policy} bounds from below.

\section{Hardness}\label{sec:hardness}

\paragraph{Exact decomposition.} For any inputs, the device output admits the exact decomposition
\begin{equation}\label{eq:decomp}
D_{ij} \;=\; \mathcal R\Big(\, S_{ij} \;-\; \sum_{g}\sum_{u\in g}\rho_u \;-\; \sum_g \gamma_g \,\Big),
\end{equation}
where $S_{ij}=\sum_u p_u$ is the exact sum, $\rho_u = p_u - \Trunc(p_u,\eta_g-W)$ is the
per-product truncation remainder in group $g$, $\gamma_g$ is the accumulator re-truncation and
per-group rounding remainder, and $\mathcal R$ denotes the per-group $\RZ$ chain. (This
reproduces the model of Section~\ref{sec:device} bit-for-bit; see Appendix~\ref{app:validation}.)
For low-rank-plus-noise operands the adversary obtains $S_{ij}$---and prefix sums of it, hence
every $\eta_g$---in $O(r)$ work per cell. The hardness therefore rests entirely on the correction
terms $\sum_u\rho_u$ and $\sum_g\gamma_g$.

\paragraph{Non-contractibility.} The decisive property is that the correction is a
\emph{per-product nonlinearity applied before the reduction}. The map $p_u\mapsto\Trunc(p_u,\eta_g-W)$
does not commute with summation:
\[
\sum_u \Trunc(p_u,\eta_g-W) \;\neq\; \Trunc\Big(\sum_u p_u,\ \eta_g-W\Big)\quad\text{in general.}
\]
Consequently $\sum_u\rho_u$ is not a function of the contracted sum $S_{ij}$, and it is not
itself a bilinear form in the operand rows. No matrix multiplication---\FP{16}, integer, or
binary---evaluates it: a contraction can produce $\sum_u p_u$ or any fixed linear image of the
products, but not the sum of a nonlinear, exponent-dependent function applied to each product.
The corrections must be computed elementwise, on the $mk$ (resp.\ $nk$) products, outside the
tensor-core datapath. This is the same obstruction that makes the integer transcript scheme
hard~\cite{KW25}; here the device applies it automatically to every MAC.

\paragraph{Assumption.} We formalize the resulting premise, in the style
of~\cite[Assumption~6.4]{KW25} and the \FP{8} transcript-unpredictability assumption.

\begin{assumption}[Accumulation-unpredictability]\label{asm:acc}
Fix $\Device=\textsf{A100}$, $\Quant$ with \FP{16} storage, noise rank $r=32$, and admissible
tile dimensions. For operands drawn as in Appendix~\ref{app:noise} (low-rank noise added before
quantization, any miner-chosen $A,B$ admitted by the policy of Section~\ref{sec:policy}), no
algorithm reconstructs the device tile $\{D_{ij}\}$ bit-exactly in expected time substantially
below the honest cost $|I_A|\,|I_B|\,k$ MACs. In particular the corrections
$\sum_u\rho_u,\sum_g\gamma_g$ of \eqref{eq:decomp} cannot be batched onto any tensor-core path.
\end{assumption}

\paragraph{Evidence.} Appendix~\ref{app:validation} summarizes the experiments supporting
Assumption~\ref{asm:acc}: (i) the exact model is validated on A100 silicon over $>10^7$ dot
products; (ii) for every input strategy tried, including sixteen adversarial constructions, the
\FP{16} breakpoint density stays in $[0.36,0.51]$ and a deliberately attacker-favorable
certified-work ratio (Appendix~\ref{app:rho}) stays $\ge 1.38$; (iii) the cheapest shortcut
predictors (exact sum rounded to \FP{32}; group sums without per-product truncation) reproduce
at most a few percent of cells; (iv) the truncation-correction attack, costed concretely, is
$6\text{--}13\times$ honest even under compute-bound-optimistic assumptions, because the
corrections run on CUDA cores rather than tensor cores. Assumption~\ref{asm:acc} is a
device-specific hardness conjecture and should be treated as such.

\section{Jackpot policy}\label{sec:policy}

Beyond the idealized setting, the miner may choose $A,B$ to make the accumulation cheaply
predictable---chiefly by arranging that few products are truncated. The verifier defends with a
pointwise policy evaluated \emph{during the bit-exact tile replay} of Section~\ref{sec:device},
when all intermediate $\eta_g$, $\bar p_u$, $\bar c$, and $S$ are available at no extra cost. The
policy approves exactly when every check below passes; it runs after the Infinity/NaN rejection.

\paragraph{Recorded during replay.} For each cell $(i,j)$ and group $g$ the verifier records:
\begin{itemize}[nosep]
\item $\textsf{bp}_{ijg}\in\{0,1\}$: whether the step is a breakpoint (Section~\ref{sec:device});
\item $\textsf{ptc}_{ijg}\in\{0,\dots,G\}$: the number of products in the group with a nonzero
      truncation remainder $\rho_u$.
\end{itemize}

\paragraph{Shared checks.} The \FP{8} entry-liveness and noise-floor checks carry over to bound
degenerate operands:
\[
\sigma_i^A\ge\sigmin\ (\forall i\in I_A),\qquad \sigma_j^B\ge\sigmin\ (\forall j\in I_B),
\qquad |\mathcal D_A|\le\eps_{\text{idle}}|I_A|k,\quad |\mathcal D_B|\le\eps_{\text{idle}}|I_B|k,
\]
with $\sigmin=1$, $\eps_{\text{idle}}=1/64$, $\tau_{\text{idle}}=8$ and the idle sets $\mathcal D_X$
as in~\cite{PearlFP8}. These keep every opened row carrying at least one quantized unit of noise
and bound the fraction of noise-dominated entries.

\paragraph{Unpredictable accumulation steps (new).} Define the tile breakpoint density and the
certified-work ratio
\[
f_{\text{bp}} := \frac{1}{|I_A||I_B|\,G^{-1}k}\sum_{i,j,g}\textsf{bp}_{ijg},
\qquad
\rho := \frac{1}{|I_A||I_B|\,k}\sum_{i,j}\Big[\,G\cdot N^{ij}_{\text{bp}} + r\cdot N^{ij}_{\text{run}} + N^{ij}_{\text{pt}}\,\Big],
\]
where $N^{ij}_{\text{bp}}$ is the number of breakpoint steps in cell $(i,j)$, $N^{ij}_{\text{run}}$
is the number of maximal runs of non-breakpoint steps, and $N^{ij}_{\text{pt}}=\sum_g\textsf{ptc}_{ijg}$
counts truncated products outside breakpoints. The ratio $\rho$ is the attacker-favorable
lower bound on reconstruction cost derived in Appendix~\ref{app:rho}. The check requires
\begin{equation}\label{eq:policyfp16}
f_{\text{bp}}\ \ge\ 0.30,\qquad \rho\ \ge\ 1.2.
\end{equation}

\paragraph{Calibration.} On A100 \FP{16}, honest and realistic workloads sit at
$f_{\text{bp}}\in[0.44,0.48]$ and $\rho\in[1.64,1.83]$, with $f_{\text{bp}}$ concentrating to
within $0.003$ across noise draws, so \eqref{eq:policyfp16} is a safety net rather than an active
constraint: every adversarial construction tested also passes it, the worst observed being
$f_{\text{bp}}=0.357$, $\rho=1.38$ (Appendix~\ref{app:validation}). Because the check is a density
gate over the whole tile, honest tiles do not fail on fluctuations, and because it concentrates,
grinding noise draws yields no advantage on honest-like data. \textsf{BF16} does not admit a
separating threshold of this form at $k\le 1024$ (Section~\ref{sec:compare}); the \FP{16}
thresholds \eqref{eq:policyfp16} apply only to the \FP{16} $\Quant$ value.

\section{Attack vectors}\label{sec:attacks}

\paragraph{Low-rank inputs.} A miner may choose low-rank $A,B$ so that $S_{ij}$ is cheap. This is
allowed; hardness does not depend on the exact sum being hard. By \eqref{eq:decomp} the device
output still requires the per-product corrections, which are non-contractible.

\paragraph{Truncation / bit-slice correction.} The corrections $\rho_u$ are shallow (median
$\sim\!4$ discarded bits per \FP{16} product), which suggests batching them on binary or int8
tensor cores via bit-planes. This fails for the structural reason of Section~\ref{sec:hardness}:
bit-planes accelerate \emph{contractions}, and the correction is not a contraction. Even granting
the adversary free exact prefix sums and a free oracle for which steps are breakpoints, the
corrections must be evaluated elementwise on CUDA cores, costing $6\text{--}13\times$ honest work
under compute-bound-optimistic accounting and far more when memory-bound
(Appendix~\ref{app:validation}).

\paragraph{Exact-sum shortcut.} Computing $\RZ_{\FP{32}}(S_{ij})$ and returning it reproduces the
device output in a negligible fraction of cells ($\le 1.6\%$ in all \FP{16} tests), because the
per-product truncation changes almost every cell.

\paragraph{Minimizing truncation.} The miner may seek operands with flat exponents or few
significant bits to reduce $f_{\text{bp}}$. For \FP{16} this is self-defeating: the mandated noise
(rank $r=32$, relative scale $\delta=1/2$) spreads product exponents enough that truncation
remains dense, and the policy \eqref{eq:policyfp16} rejects any tile that nonetheless falls into a
low-truncation regime. For \textsf{BF16} the same strategy succeeds, which is why \textsf{BF16} is
excluded (Section~\ref{sec:compare}).

\paragraph{Device emulation and ASICs.} Reproducing A100 accumulation on another GPU requires
emulating the $G=8$ grouping and per-product truncation, which that device's tensor cores cannot
do cheaply; it must fall back to the same elementwise path as any attacker. A custom ASIC could
implement the per-product correction more cheaply than a full \FP{16} multiplier, so the argument
is against commodity GPUs, not arbitrary silicon; this is the same economic question raised
for the \FP{8} scheme and in the Pearl open-problems note.

\paragraph{Seed and commitment grinding.} The noise is seed-derived after commitment; grinding
seeds only yields additional lottery tickets, bounded by the usual Fiat--Shamir argument.

\section{Honest miner overview}\label{sec:miner}

The intended miner mines with \FP{16} multiplications it wants computed anyway.
\begin{enumerate}[nosep]
\item \textbf{Prequantize.} Encode $A,B$ in the \FP{16} commit format (Appendix~\ref{app:commit}).
\item \textbf{Commit.} Hash $B$ under a recent header and derive $\texttt{noise\_seed}_B$; commit
      $A$ under the proposed header and derive $\texttt{noise\_seed}_A$.
\item \textbf{Noise and quantize.} Derive the rank-$r$ noise and compute
      $\At=\NoisyQuantize(A,\dots)$, $\Bt=\NoisyQuantize(B,\dots)$ in $O((m+n)kr)$ MACs
      (Appendix~\ref{app:noise}).
\item \textbf{Multiply.} Run the A100 \FP{16} kernel once on the noised operands---$mnk$ MACs, the
      dominant cost, identical to the kernel the job would run without the protocol.
\item \textbf{Scan for winners.} Cut the output into tiles, compute $\Extract$ and the jackpot
      hash per tile, and compare against the scaled target.
\item \textbf{Open on a win.} For a winning tile admitted by the policy, assemble the proof.
\end{enumerate}
Recovering the useful product costs the low-rank correction plus the $O((mn+mk+nk)r)$ overhead of
hashing, noising and scanning, for a $1+o(1)$ total overhead when $m,n,k\gg r$. The recovered
product carries \FP{16}-level accuracy: the device truncation that secures the tile is of the same
order as---and is reconstructed together with---the quantization error, so no accuracy beyond
ordinary \FP{16} rounding is sacrificed. Critically, the miner's kernel must use the committed
instruction order with no split-$k$ or atomic accumulation, since any reordering changes the
bit-exact output (Appendix~\ref{app:arith}).

\section{FP16 vs.\ FP8 vs.\ BF16}\label{sec:compare}

\paragraph{Why \FP{16} works.} \FP{16} products carry $22$ significant bits against the $24$-bit
accumulator window, so truncation is dense and, crucially, insensitive to the miner's choices:
across all tested strategies $f_{\text{bp}}\in[0.36,0.51]$. The non-contractibility argument then
applies with a wide margin.

\paragraph{Why \textsf{BF16} does not.} \textsf{BF16} products carry only $16$ bits and fit inside
the $24$-bit window unless more than $8$ binades below $\eta$. Flat rank-one sign-pattern operands
leave $>\!99\%$ of products exact and only $6\text{--}8\%$ of steps as breakpoints; a groupwise
model without per-product truncation already reproduces more than half the outputs, and the
certified-work ratio falls to $\rho\approx 0.30$. These operands pass the liveness and noise-floor
checks. Honest \textsf{BF16} is itself marginal ($\rho\in[0.84,1.10]$ at $k=1024$), so no
threshold separates honest from adversarial \textsf{BF16} with margin. \textsf{BF16} mining, if
desired, requires either \FP{16} storage, $k\ge 4096$, or a larger noise rank; this document does
not specify it.

\paragraph{Relation to \FP{8}.} The \FP{8} scheme remains the cheapest option and its hardness
margin is ample; it is preferred for workloads tolerant of \FP{8}-residual accuracy. \FP{16}
trades a modest constant factor (denser truncation to reconstruct, a slightly larger proof) for
\FP{16}-level output accuracy and a hardness source independent of the rounding grid. A deployment
may offer both $\Quant$ values and weight their targets so that each is mined at its true device
throughput; the economics mirror the \FP{8} discussion in~\cite{PearlFP8}.

\appendix

\section{A100 arithmetic specification}\label{app:arith}
Element-wise \FP{16}/\textsf{BF16}/\FP{32} operations follow IEEE-754-2019 with the roundings
used by the device: casts to \FP{16} round to nearest even and saturate per the format; the
matrix-multiply accumulation is specified in Section~\ref{sec:device}. The verifier rejects NaN
and infinity in operands and in every intermediate, applying the test after rounding. The A100
accumulator keeps $W=24$ significant bits (no sub-\FP{32} truncation at the accumulator level,
unlike the Hopper/Blackwell \FP{8} paths of~\cite{PearlFP8}); the nonlinearity is entirely in the
per-product, per-group truncation of Section~\ref{sec:device}. The \texttt{m16n8k8} instruction is
identical to \texttt{m16n8k16} with $G=8$; the \FP{16}-accumulate variant (\texttt{HMMA.16816.F16})
uses the same grouping and truncation followed by round-to-nearest-even to \FP{16} per group and is
not used by this scheme.

\section{FP16 commit type}\label{app:commit}
An operand with $n$ rows is committed as a single matrix $V\in\FP{16}^{n\times k}$, each row a
Merkle leaf-set under the key and hash identifier of the main protocol, with aggregate root $H_X$.
Unlike the \FP{8} FP10 commit (int8 values plus a block scale), the \FP{16} format stores values
directly, since \FP{16} already carries the target precision; no block scale is required. Row norms
$\ell_2(X_t),\ell_\infty(X_t)$, used by the noisy-quantization scales, are computed from $V$ in
\FP{32} and floored as in~\cite{PearlFP8}.

\section{Noise and noisy quantization}\label{app:noise}
The noise sampler, seed chain, and rank-$r$ construction $N=EF^\T$ are those of the \FP{8}
specification, with $r=32$ and per-row relative scale $\delta=1/2$. Noisy quantization computes,
per row, $\At_t=\RNE_{\FP{16}}(\alpha_t X_t+\beta_t N_t)$ with the scales $\alpha_t,\beta_t$ chosen
so the noise carries relative Euclidean weight $\delta$; the clamp and norm-floor rules follow
Appendix~C.4 of~\cite{PearlFP8} with target magnitude $Q$ set to the largest finite \FP{16} value
$65504$. Because \FP{16} is the target precision, $\alpha_t$ performs only range normalization and
the recovered product retains \FP{16} accuracy after the low-rank correction.

\section{Certified-work ratio}\label{app:rho}
The ratio $\rho$ of Section~\ref{sec:policy} is an attacker-favorable lower bound on the cost of
reconstructing a tile under the following generous premises: (i) exact prefix sums over any
$k$-range cost $r$ MAC-equivalents per cell via the low-rank structure; (ii) a run of consecutive
non-breakpoint steps can be merged at cost $r$ plus one per truncated product; (iii) each
breakpoint step costs $\min(G,r)=G$ to resolve its discarded low bits from the previous exact
state; (iv) the breakpoint locations are supplied for free by an oracle. Summing gives the
per-cell cost $G\,N_{\text{bp}}+r\,N_{\text{run}}+N_{\text{pt}}$ used in the definition, normalized
by the honest $k$. Premises (i)--(iv) overstate the attacker's power---in particular, quantization
makes the operands only approximately low rank---so a measured $\rho\ge 1.2$ is a conservative
certificate that no such shortcut beats honest work.

\section{Empirical validation}\label{app:validation}
\paragraph{Scope of the bit-exactness claim.} The accumulation model is validated bit-exact
\emph{only} on GA100 (\texttt{sm\_80}) silicon. Other Ampere/Ada parts (\texttt{sm\_86},
\texttt{sm\_87}, \texttt{sm\_89}) are distinct tensor-core designs whose intra-instruction
accumulation has not been characterized against this model---they may or may not agree---so the
whitelisted device is GA100 (\texttt{sm\_80}) exactly, and the miner admits that capability alone
(fail-closed). Whitelisting any further SKU requires repeating the validation below on that part.

All experiments were run on A100-class (GA100, sm\_80) silicon. The accumulation model of
Section~\ref{sec:device} matched the hardware bit-for-bit on $>1.1\times10^6$ random dot products
per device (both devices agreeing), on $69$ directed edge-case scenarios totaling $6.9\times10^6$
dot products (subnormal inputs and accumulators, exact cancellation, mixed signs, grid-boundary
products, partial and empty groups, $k=2048$ chains, overflow, infinity, NaN, and $0\cdot\infty$),
and on the adversarial and realistic noised operands below. One residual is known: \textsf{BF16}
chains whose partial sums remain entirely in the \FP{32}-subnormal range can be off by one
subnormal ulp in $<\!1\%$ of outputs; this cannot occur for \FP{16} or for $O(1)$-scaled operands
and is excluded by policy.

Adversarial study (tiles $16\times16$, $k=1024$, $r=32$, $\delta=1/2$, operands row-scaled to unit
max magnitude): sixteen input strategies---realistic Gaussian with outlier channels, heavy-tailed,
$A=B=0$, rank-two, several flat and few-bit constructions designed to minimize truncation, sparse,
spiked, and geometric---all kept \FP{16} at $f_{\text{bp}}\in[0.36,0.51]$ and $\rho\ge 1.38$, with
the cheapest shortcut predictors reproducing $\le 1.6\%$ of cells (one geometric case $13\%$, still
$\rho=1.68$). The same strategies drive \textsf{BF16} to $f_{\text{bp}}$ as low as $0.06$ and
$\rho\approx 0.30$. The truncation-correction attack, costed end-to-end, is $6\text{--}13\times$
honest (compute-bound-optimistic) and $220\text{--}470\times$ (measured, memory-bound), because the
per-product corrections cannot use tensor cores.

\paragraph{Reproducible harness.} A self-contained harness is committed alongside this
specification at \texttt{docs/fp16\_scheme/validation/} (it links only the CUDA runtime, so it runs
on an \texttt{sm\_80} box without the miner's Python stack). On \textsf{A100}-class silicon
(validated on a GA100 \texttt{sm\_80} device) it:
(i) captures the output of the real \texttt{mma.sync.m16n8k16.f32.f16} and checks the model of
Section~\ref{sec:device} against the device \emph{bit-for-bit} over $910$ dot products---the $238$
committed reference vectors plus $272$ generated edge cases ($k$ up to $1024$, subnormal operands,
subnormal accumulators, subnormal outputs, near-total cancellation, grid-boundary, and
max-magnitude)---with zero mismatches, so the model is validated against the device rather than
against another copy of itself;
(ii) resolves the rounding mode on silicon, \emph{confirming round-toward-zero} at both the
per-product alignment and the per-group \FP{32} encode: the device matches the \RZ model on all
$400$ vectors constructed to separate \RZ from round-to-nearest-even, while each RNE variant
disagrees with the device on a large fraction of them;
(iii) recomputes the policy metrics $f_{\text{bp}},\rho$ and the shortcut-predictor rates over the
adversarial strategies, through both an independent noise implementation and the exact consensus
noise pipeline, reproducing the regime above ($f_{\text{bp}}\in[0.45,0.52]$,
$\rho\in[1.67,1.94]$, cheapest shortcut $\le 1.6\%$); and
(iv) micro-benchmarks the truncation-correction attack against the honest tensor-core \FP{16} GEMM
(the miner's own kernel), measuring $\approx 38\text{--}43\times$ honest on this silicon even when
the attacker is granted the exact sums and the breakpoint oracle for free---consistent with the
accounting of Section~\ref{sec:attacks}. The model-vs-silicon corpus is enforced in CI by the Rust
test \texttt{api::fp16::accumulate::matches\_hardware\_capture\_k1024}; the overflow guard was never
triggered, confirming that a single \FP{16} dot product cannot overflow.

\begin{thebibliography}{9}
\bibitem{KW25} I.~Komargodski and O.~Weinstein. \emph{Proofs of Useful Work from Arbitrary Matrix
Multiplication}. arXiv:2504.09971; IACR ePrint 2025/685, 2025.
\bibitem{PearlFP8} Pearl Research Team. \emph{Pearl Floating Point Scheme Specification}
(FP8 certificate-v4), September 2026.
\bibitem{Hawkeye} E.~Badash, D.~Boneh, I.~Komargodski, and M.~Srivastava. \emph{Hawkeye:
Reproducing GPU-level Non-determinism}. arXiv:2603.20421, 2026.
\bibitem{BLAKE3} J.~O'Connor, J.-P.~Aumasson, S.~Neves, and Z.~Wilcox-O'Hearn. \emph{BLAKE3:
one function, fast everywhere}, 2020.
\end{thebibliography}

\end{document}
